\chapter{Discussion}
\label{chap:discussion}

The objective of this thesis was to find an effective yet cheap recipe for
turning a pretrained vision-language model into an acting agent, using Minecraft
as a deliberately hard test bed. The recipe was held to two constraints: the
backbone stays frozen, and training stays inexpensive enough to sweep the full
ablation on a single rented GPU (Section~\ref{caching}). This chapter reads the
results of Chapter~\ref{chap:experiments} against that objective, asks why the
two backbones behave so differently, examines what the agents do and do not
achieve in the chop environment, states the limits of the study, and lays out
two directions the findings open up.

\section{CLIP and LLaVA Head-to-Head Under a Cheap, Frozen Regime}
\label{sec:disc-headtohead}

The clearest message of the ablation is that, under the constraints this thesis
imposes, the simpler backbone wins --- with the important qualification,
developed below, that this is a property of the minimal head, not an intrinsic
ranking of the two backbones. Of the two frozen extractors, the
contrastive CLIP encoder both conditions on the prompt
(Section~\ref{sec:exp-cond}) and collects the discriminating item
(Section~\ref{sec:exp-task}), whereas the joint LLaVA-1.5 backbone does neither
in its matched, knob-free configuration. On the task that separates the cells,
dirt collection, CLIP digs reliably with and without the prompt ($2.53$ and
$1.60$ per episode, averaged over three training seeds) while both LLaVA anchor
cells sit at the floor ($0.07$, $0.47$); on conditioning, CLIP with language
shifts its \texttt{attack} rate by $42$ percentage points between the empty and
on-task prompts, an effect an order of magnitude larger than anything either
LLaVA cell produces and present in all three seeds. For a practitioner who wants
a controllable short-horizon agent on a small budget, the result is encouraging:
the cheaper $0.43$B contrastive encoder, paired with a $4.6$M-parameter head, is
the more effective choice here \emph{under this minimal head}, not the $7$B
autoregressive one. Set against the contractor it imitates, though, even the best
agent is modest in absolute terms: measured from the recordings, the demonstrator
collects $18.3$ dirt per thousand steps, so CLIP\,+\,lang's $2.5$ is about a
seventh of the expert rate --- far above the random floor of zero, but a reminder
that the comparison this study makes cleanly is between cells, not against
mastery.

This separation is invisible to held-out frame accuracy. On the test frames the
four cells are near-identical --- movement-$F_1$ spans only $0.368$ to $0.431$ and
even ranks CLIP $-$ lang first --- yet they diverge from a dug stack of dirt to
nothing in rollout (Section~\ref{sec:exp-offline}). The head-to-head verdict is therefore a
property of closed-loop behaviour, not of frame-level fit, and it would be missed
entirely by the offline metrics that imitation-learning studies usually report.
This is itself the answer to the third research question, and a caution against
selecting backbones on frame accuracy alone.

This conclusion must be stated at the right altitude, and the FiLM probe of
Section~\ref{sec:exp-film} fixes that altitude. The probe shows that the same
frozen LLaVA features that produced a hard zero can be made to collect dirt at
CLIP-comparable levels once the head reads them through feature-wise modulation
rather than plain concatenation. The honest claim is therefore not that LLaVA's
representation is intrinsically weaker for control, but that \emph{under the
simplest head---the one the controlled $2\times2$ holds fixed---CLIP extracts
more usable task signal per unit of engineering effort}. The advantage CLIP
shows is an advantage of the cheap regime, which is precisely the regime this
thesis set out to characterise. A contrastive encoder aligns image and text in a
single shared embedding by construction, so a linear head can read a
prompt-dependent direction out of it directly; the joint decoder hides the same
information behind a pooling and read-out step that the minimal head cannot
invert.

\section{Why the LLaVA Language Pathway Stays Silent}
\label{sec:disc-washout}

The most striking negative result is that LLaVA's behaviour is statistically
indistinguishable across all four prompt conditions: the explicit task string
moves the policy no more than the zeroed control does. The tempting reading ---
that the language tokens are drowned out by the roughly $576$ image tokens --- is
too crude. The \texttt{<image>} placeholder is expanded into those visual tokens
\emph{alongside} the prompt tokens, the split-pooling of Section~\ref{sec:backbone}
keeps the text-role tokens as an explicit half of the feature, and the linear
probe of Section~\ref{sec:exp-probe} shows the prompt is in fact almost perfectly
recoverable from that half ($99.7\%$, collapsing to chance when it is zeroed). The
prompt is not washed out in the crude sense.

Recoverable, however, is not the same as usable, and the rest of the evidence
shows the limit is the representation, not the head. In the combined training set
the two prompts are perfectly correlated with the two tasks' imagery
(tree-felling versus dirt-digging scenes), and the pooled LLaVA text half is
image-entangled --- its tokens attend to the frame inside the decoder --- so the
$99.7\%$ is largely a prompt--image confound rather than a clean prompt axis. The
rollout evaluation breaks that correlation: it issues all four prompts in one
fixed scene (Section~\ref{sec:exp-setup}), so prompt and image are decorrelated,
and the LLaVA policy still does not move --- the prompt axis is absent there, not
merely masked by the image. Two converging checks confirm the limit is upstream
of the head rather than in it. First, a FiLM head that multiplicatively gates the
features on that very text vector, strictly more expressive than the concatenating
baseline, still does not condition (Section~\ref{sec:exp-film}); if a usable
prompt axis were present a gating head would exploit it. Second, that same gating
head, and a second alternative read-out, \emph{do} recover the dirt \emph{task}
from the identical features --- so the heads are capable enough to use what the
representation exposes, and what it exposes is task identity, not a steerable
prompt direction. Head capacity is therefore not the limiting factor. The contrast
with CLIP isolates what is: CLIP's text encoder is physically separate from its
image encoder, so its prompt embedding is an image-independent axis that moves
when, and only when, the prompt changes, and the head learns to ride it. LLaVA's
final-layer text positions, mean-pooled, fold the prompt and the frame together,
so changing only the prompt at a fixed scene barely moves the representation. The
bottleneck is thus upstream of the head, and the remedies that follow are
representational rather than architectural on the head side: pooling the prompt
from a position that is not image-saturated --- the final text token rather than
the mean --- or lightly fine-tuning the backbone so the text pathway carries a
separable signal. A counterfactual probe that swaps the prompt at a fixed frame
measures this directly (Section~\ref{sec:exp-probe}): LLaVA's pooled text half
barely moves under the swap (cosine $1.00$) while CLIP's moves substantially
(cosine $0.74$), confirming that the limit is the representation, not the head.

\section{Targeting Without Completion in the Chop Environment}
\label{sec:disc-chop}

Reward and item counts tell only part of the story for the chop task, where every
cell is at the floor (Section~\ref{sec:exp-task}). Inspecting the rollout videos
rather than the scalar logs reveals that the better CLIP agents are not failing
at random: they orient toward nearby trees, walk up to a trunk, and fire
\texttt{attack} into it --- the coarse structure of chopping is present. What they
miss is the fine control that finishes the job: holding the camera steady on a
single block long enough, at the right pitch, for the break to complete. The
agent targets the tree but does not reliably mine the block. This is the
signature of the compounding-error ceiling that behavioural cloning carries by
construction (Section~\ref{sec:bg-il}): trained only to imitate single-step
actions with no reward and no corrective feedback, the policy reproduces the
gross demonstrated behaviour but drifts on the tight closed-loop control that
block-breaking demands, and nothing in the offline objective pushes it back onto
the narrow trajectory that completes the chop. A demo-side analysis locates where
that trajectory runs out of supervision. Aligning camera motion to the onset of
sustained \texttt{attack} runs --- the median chop commit lasts about $90$ frames
--- shows the contractor does move the camera more in the half-second \emph{before}
committing to a chop than during it ($3.1^\circ$ versus $1.4^\circ$ of per-frame
motion), so an aiming phase is present in the data. But that motion is largely
undirected: its net signed pitch is only $-0.45^\circ$ against $1.4^\circ$ of
magnitude --- a faint downward bias toward the trunk base, swamped by exploratory
wander --- and yaw carries no consistent direction at all. The precise
\enquote{aim at the trunk} signal a head would need to finish the block is
therefore weak in the demonstrations themselves, so part of the chop gap is a
limit of the supervision, not only of the cloned policy. That the qualitative
behaviour is already on the right track, while the scalar reward stays at zero, is itself a
finding: it suggests the gap to task completion here is one of refinement, not of
gross capability, and that the right lever is a training signal that can reward
the last few degrees of aim rather than a larger or unfrozen backbone.

The conditioning result of Section~\ref{sec:exp-cond} casts this task in a second,
initially puzzling light: the on-task \enquote{chop a tree} prompt \emph{suppresses}
the CLIP\,+\,lang agent's \texttt{attack} key, from $80.4\%$ under the empty prompt
to $38.2\%$ --- a drop in precisely the action that chopping requires. A natural
hope is that this merely replays the demonstrations, but the contractor data rules
that out. Measured directly from the recordings, the chop demonstrator attacks on
$83.1\%$ of steps and the dirt demonstrator on $94.6\%$; both tasks are
attack-heavy, and the empty-prompt policy ($80.4\%$) faithfully reproduces that
baseline. The prompt then drives \texttt{attack} \emph{below} either demonstrator,
so the suppression cannot be imitation of a low-attack chopping style --- no such
style exists in the data. It is a genuine prompt-induced reorganisation of
behaviour rather than a replayed statistic; what it is not is evidence of
\emph{correct} task behaviour, since neither the high-attack empty-prompt agent nor
the suppressed prompted one completes a chop. The $42$-point shift thus shows that
the prompt steers the policy (RQ2) without showing that it steers it toward the
goal: the chop environment exposes the limit of conditioning as sharply as the
dirt environment exposes its use, and steering and task-appropriateness remain
separate properties.

\section{Limitations}
\label{sec:disc-limits}

Several limits bound how far these conclusions generalise. First, the
in-environment comparisons rest on ten episodes per cell per training seed;
although MineRL is seed-deterministic, so each run reproduces exactly, ten
episodes is still a small sample over ten distinct world seeds. The per-step
firing rates that carry the conditioning result are estimated from
$\sim\!3\times10^4$ pooled observations and are robust, but the per-episode task
metrics are noisier, so the within-CLIP prompt-versus-no-prompt difference is only
suggestive and the cross-backbone task gap is the one task comparison I treat as
reliable. To bound the training-side counterpart of this caveat, each cell is
trained at three independent seeds (head initialisation and data-order shuffle);
the headline effects survive it. The conditioning shift is $-42.3\pm10.4$ points
and large in every seed ($h=-1.04,-1.12,-0.57$); the LLaVA flats are symmetric
across all three seeds ($|h|\le0.16$); and the cross-backbone dirt gap holds (CLIP
$2.53\pm1.25$ and $1.60\pm0.50$ versus the LLaVA anchors' $0.07$ and $0.47$). What
three seeds reveal, and a single seed hid, is that several absolute numbers were
seed-lucky: CLIP\,+\,lang's dirt count was $4.3$ in the original seed but
$2.53\pm1.25$ across seeds, and the FiLM task recovery likewise. The signs and
orderings are stable; the precise magnitudes should be read with their across-seed
spread. The quantitative \emph{task} story, in turn, rests on a single working
task: chop is at the floor for every cell, so dirt collection is the only task on
which the agents produce a completion metric at all, with the chop environment
contributing the per-step conditioning shift and the qualitative targeting result
but no scalar separation. Second, the study fixes the
backbone as frozen and the head as a minimal two-layer MLP; it therefore
characterises the cheap-and-frozen corner of the design space and cannot speak to
what either backbone achieves under fine-tuning or a richer head, a boundary the
FiLM probe already shows matters. Third, the policy is single-frame with a short
past-action window, so it is structurally ill-suited to the long-horizon credit
assignment that full task completion in Minecraft demands. Fourth, the two tasks
are short-horizon and the prompt doubles as a task-identity signal on the combined
dataset (Section~\ref{ablation}), so the conditioning result demonstrates that
language \emph{routes behaviour by task} for CLIP, not that the model parses
compositional instructions. Relatedly, the dirt-collection metric is not a pure
measure of goal pursuit: every cell digs some dirt incidentally, by attacking the
ground, even under the empty prompt and in the same fixed world
(Table~\ref{tab:task}), so an absolute dirt count conflates goal-directed digging
with incidental ground-attacking. The prompt-driven component is visible only for
CLIP\,+\,lang, and even there it is weak and seed-variable at the episode level
(an empty-to-dirt-prompt lift of $1.4\pm1.7$), which is why the reproducible
prompt-steering evidence is the per-step firing-rate conditioning, not the item
count. Finally, the LLaVA no-prompt control is not a perfect
negative: because the prompt still flows through the decoder and only the pooled
text vector is zeroed, the LLaVA image features remain prompt-conditioned even in
the control, which is one more reason the cleanest statement is the
\emph{symmetric} flatness of LLaVA across all four conditions rather than a single
prompt-versus-control contrast.

\section{Future Directions}
\label{sec:disc-future}

The findings push in two complementary directions.

\paragraph{Earn the joint backbone its keep.} The silence of LLaVA's language
pathway is a property of the frozen, minimal-head regime, not a verdict on joint
vision-language models, which have driven strong long-horizon Minecraft and
robotics agents once a heavier pipeline is allowed
\cite{li2025jarvisvla, kim2024openvla, brohan2023rt2, black2024pi0}. Because the
bottleneck this study locates is upstream of the head, in the pooled text
representation (Section~\ref{sec:disc-washout}), the cheapest next step keeps the
whole two-stage pipeline and changes only how the prompt is pooled: re-caching from
a position that is not image-saturated --- the final text token rather than the
mean over text positions --- might expose the prompt as a separable axis at no
extra training cost. Beyond that, relaxing the freeze with parameter-efficient
fine-tuning such as LoRA \cite{hu2022lora}, or a few learnable conditioning tokens,
could let the decoder carry a behaviourally separable prompt signal without the
feature distortion that full fine-tuning risks
\cite{kumar2022finetuningdistortpretrainedfeatures, simchowitz2025pitfalls}.
These remedies are now well motivated rather than speculative, because the
counterfactual prompt-swap probe of Section~\ref{sec:exp-probe} has confirmed
their target directly: holding the frame fixed and swapping only the prompt moves
LLaVA's pooled text half by a mere one percent (cosine $1.00$), against CLIP's
substantial swing (cosine $0.74$), so the prompt does not register in LLaVA's
pooled representation at all. The fix must therefore act on that representation ---
re-pooling from the final text token or relaxing the freeze so the text pathway
carries a prompt-separable signal --- and the same probe gives a cheap acceptance
test for whether it worked: a successful intervention should raise the LLaVA
prompt-swap displacement toward CLIP's before any head is retrained.

\paragraph{Lean into cheap models, and compose them.} The other direction takes
the CLIP result at face value: a small contrastive backbone with a light head is
already an effective short-horizon controller, so it is worth pushing how far that
cheapness scales. Long-horizon competence need not come from a single large model;
it can come from composing several cheap, prompt-conditioned CLIP heads into a
multi-skill policy that is directed by a separate language model acting as a
high-level planner \cite{lifshitz2023steve1, fan2022minedojo, lynch2021langlfp}.
Such a hierarchy keeps each control policy simple and inexpensive to train while
letting the planner supply the knowledge-based reasoning that bridges abstract
goals and low-level actions --- exactly the gap a frozen contrastive encoder
cannot close on its own.

\paragraph{Close the loop with reward.} Cutting across both directions is the
chop observation of Section~\ref{sec:disc-chop}: a behaviourally-cloned CLIP agent
that already targets trees but cannot finish the block is a strong starting point
for a second, online training stage. Fine-tuning the cloned policy against a
rollout-based reward, or nudging it with corrective feedback in the
DAgger family \cite{ross2011}, would supply exactly the closed-loop signal that
offline cloning lacks, reinforcing the last few degrees of aim that separate
visible targeting from a completed chop. Because the underlying gross behaviour is
already present, such a stage plausibly needs far less interaction than learning
the task from scratch --- a hypothesis that follows directly from the qualitative
result and that the cheap-and-frozen recipe of this thesis is well positioned to
seed.
